\documentclass[11pt]{article}
\usepackage[a4paper,margin=25mm]{geometry}
\usepackage{amsmath,amssymb,amsthm}
\usepackage[hidelinks]{hyperref}
\hypersetup{pdftitle={H34\_01: y\textasciicircum 3+y\textasciicircum 2=x\textasciicircum 4+x+4 has no integer solutions}}
\newtheorem{theorem}{Theorem}
\newtheorem{lemma}[theorem]{Lemma}
\newcommand{\Z}{\mathbb{Z}}
\title{\textbf{H34\_01}\\[4pt] The equation $y^3+y^2=x^4+x+4$ has no integer solutions}
\author{}
\date{}
\begin{document}
\maketitle
\vspace{-8ex}
\begin{center}
Size $h=34$, length $l=12$, identifier E1179 (with $y$ replaced by $-y$)
\end{center}

\begin{theorem}\label{thm:main}
The equation
\begin{equation}\label{eq:main}
y^3+y^2=x^4+x+4
\end{equation}
has no integer solutions.
\end{theorem}

The proof uses a polynomial identity of the form $HT=U^2+7V^2$ on the curve \eqref{eq:main}, and the fact that a prime $p$ with $p\equiv3$, $5$ or $6\pmod7$ divides $U^2+7V^2$ only to an even power. Apart from polynomial identities, which can be checked by expanding both sides, and finite checks of congruences, the proof uses only the law of quadratic reciprocity.

\section{\texorpdfstring{The norm form $a^2+7b^2$}{The norm form}}

For a prime $p$ and a non-zero integer $N$, we write $v_p(N)$ for the exponent of $p$ in the prime factorization of $N$.

\begin{lemma}\label{lem:divides}
Let $p$ be a prime with $p\equiv3$, $5$ or $6\pmod7$, and let $a,b\in\Z$. If $p\mid a^2+7b^2$, then $p\mid a$ and $p\mid b$.
\end{lemma}

\begin{proof}
Note that $p$ is odd and $p\neq7$. Suppose that $p\nmid b$. Choose $c\in\Z$ with $bc\equiv1\pmod p$ and put $r=ac$. Then $r^2\equiv-7(bc)^2\equiv-7\pmod p$, so $-7$ is a square modulo $p$, that is, $\bigl(\frac{-7}{p}\bigr)=1$ for the Legendre symbol. By the law of quadratic reciprocity,
\[
\Bigl(\frac{7}{p}\Bigr)=(-1)^{\frac{p-1}2\cdot\frac{7-1}2}\Bigl(\frac{p}{7}\Bigr)=(-1)^{\frac{p-1}2}\Bigl(\frac{p}{7}\Bigr),
\qquad
\Bigl(\frac{-1}{p}\Bigr)=(-1)^{\frac{p-1}2},
\]
hence $\bigl(\frac{-7}p\bigr)=\bigl(\frac{-1}p\bigr)\bigl(\frac7p\bigr)=\bigl(\frac p7\bigr)$. But $p\equiv3$, $5$ or $6\pmod7$ is not a square modulo $7$, so $\bigl(\frac p7\bigr)=-1$, a contradiction. Hence $p\mid b$, and then $p\mid a^2$, so $p\mid a$.
\end{proof}

\begin{lemma}\label{lem:even}
Let $p$ be a prime with $p\equiv3$, $5$ or $6\pmod7$, and let $a,b\in\Z$ be not both zero. Then $v_p(a^2+7b^2)$ is even.
\end{lemma}

\begin{proof}
Let $p^k$ be the largest power of $p$ dividing both $a$ and $b$, and write $a=p^ka'$ and $b=p^kb'$. Then $p$ does not divide both $a'$ and $b'$, so $p\nmid a'^2+7b'^2$ by Lemma~\ref{lem:divides}. Since $a^2+7b^2=p^{2k}(a'^2+7b'^2)$, we get $v_p(a^2+7b^2)=2k$.
\end{proof}

\section{The auxiliary polynomials}

Put $F(x,y)=y^3+y^2-x^4-x-4$, so that \eqref{eq:main} is the equation $F(x,y)=0$. Define
\begin{align*}
H&=2x+y+3,\\
T&=4y^2-8(x+1)y+16x^2+40x+24,\\
U&=2x^2+8x+9,\qquad V=2x+1,\\
G&=4,
\end{align*}
and
\[
A=2,\qquad B=-(2x+7).
\]
Expanding both sides, one checks the identities
\begin{align}
HT-U^2-7V^2&=GF,\label{eq:norm}\\
T&=4(y-x-1)^2+4(x+1)(3x+5),\label{eq:sign}\\
AU+BV&=11.\label{eq:bezout}
\end{align}

\section{Proof of Theorem~\ref{thm:main}}

Suppose that $(x,y)\in\Z^2$ is a solution of \eqref{eq:main}. Then $F(x,y)=0$, and \eqref{eq:norm} gives
\begin{equation}\label{eq:HT}
H\,T=U^2+7V^2,
\end{equation}
where $H$, $T$, $U$ and $V$ now denote the values of these polynomials at $(x,y)$.

\paragraph{Step 1: congruences.} Checking the $49$ residue pairs modulo $7$, we find that the only solutions of \eqref{eq:main} modulo $7$ are $(x,y)\equiv(2,3)$ and $(x,y)\equiv(4,2)$. In these cases $H\equiv3$ and $H\equiv6\pmod7$, respectively.

\paragraph{Step 2: signs.} For $x=-1$, \eqref{eq:main} reads $y^3+y^2=4$, and the values of $y^3+y^2-4$ at $y=0,1,2$ are $2,1,2$ modulo $3$, so $x\neq-1$. Hence the factors $x+1$ and $3x+5$ have the same sign, so $(x+1)(3x+5)>0$, and \eqref{eq:sign} gives $T>0$. By \eqref{eq:bezout}, $U$ and $V$ are not both zero, so $U^2+7V^2>0$, and \eqref{eq:HT} gives $H>0$.

\paragraph{Step 3: primes with $p\equiv3$, $5$ or $6\pmod7$.} Let $p$ be a prime with $p\equiv3$, $5$ or $6\pmod7$. We show that $v_p(H)$ is even. By \eqref{eq:bezout}, $U$ and $V$ are not both zero, so \eqref{eq:HT} and Lemma~\ref{lem:even} show that $v_p(HT)=v_p(U^2+7V^2)$ is even. Suppose that $p$ divides both $H$ and $T$. Then $p\mid U^2+7V^2$ by \eqref{eq:HT}, so $p\mid U$ and $p\mid V$ by Lemma~\ref{lem:divides}, and $p\mid 11$ by \eqref{eq:bezout}. But $11\equiv4\pmod7$, a contradiction. Hence $p$ does not divide both $H$ and $T$. If $p\nmid T$, then $v_p(H)=v_p(HT)$ is even, and if $p\nmid H$, then $v_p(H)=0$.

\paragraph{Step 4: contradiction.} By Steps~1 and~2, $H$ is a positive integer not divisible by $7$. Every prime factor $p$ of $H$ with $p\not\equiv3,5,6\pmod7$ satisfies $p\equiv1$, $2$ or $4\pmod7$, and every prime factor with $p\equiv3,5,6\pmod7$ occurs with an even exponent by Step~3. Since the set $\{1,2,4\}$ of residues modulo $7$ is closed under multiplication and contains the squares of all integers not divisible by $7$, $H\equiv1$, $2$ or $4\pmod7$. This contradicts Step~1, and the theorem is proved.
\qed

\end{document}
